\documentclass[11pt]{article}

\usepackage[margin=1in]{geometry}
\usepackage{amsmath,amssymb,amsthm,mathtools}
\usepackage[hidelinks]{hyperref}

\newtheorem{theorem}{Theorem}[section]
\newtheorem{proposition}[theorem]{Proposition}
\newtheorem{lemma}[theorem]{Lemma}
\newtheorem{corollary}[theorem]{Corollary}
\theoremstyle{definition}
\newtheorem{definition}[theorem]{Definition}
\theoremstyle{remark}
\newtheorem{remark}[theorem]{Remark}

\newcommand{\F}{F}
\newcommand{\On}{\mathrm{On}_2}
\newcommand{\rad}{\operatorname{rad}}
\newcommand{\rank}{\operatorname{rank}}
\newcommand{\Cl}{\operatorname{Cl}}
\newcommand{\BW}{\operatorname{BW}}
\newcommand{\Wq}{W_q}
\newcommand{\wpmap}{\wp}

\title{The Arf Switches Off:\\
Quadratic Forms over the Full Nimber Field}
\author{a9lim and Codex}
\date{7 August 2026}

\begin{document}
\maketitle

\begin{abstract}
This note resolves the \texttt{off} problem in Ogdoad's open ledger.  The
ground field, rather than a new transfinite trace, decides the question.
Conway's full nimber field $\On$ is algebraically closed of characteristic
$2$.  Consequently the Artin--Schreier map is surjective, every regular
finite-dimensional quadratic form over $\On$ is hyperbolic, and the
quadratic Witt group is zero.  A possibly singular form is classified by
the rank of its polar form, the radical dimension, and whether its
restriction to the radical vanishes; there is no further Arf coordinate.
The Clifford/Brauer--Wall class of every regular form is therefore split.

The finite-field Arf bit remains correct as a \emph{relative} invariant over
the chosen finite nim-subfield.  It dies after a quadratic scalar extension,
so its directed colimit through the finite nim-subfields is zero.  This also
exposes an empty mathematical middle case: every ordinal below
$\omega^{\omega^\omega}$ is algebraic over $\mathbb F_2$, so every finite
metric in that segment has a common finite field of definition.  The backend
detects that field only when the required Kummer excess rows are verified and
the checked degree arithmetic remains representable.
Finally, Ogdoad's partial \texttt{Ordinal} arithmetic need not pretend to
construct Artin--Schreier roots that lie outside its verified window, even
though the full-$\On$ classification is mathematically complete.
\end{abstract}

\section{Question and answer}

Let $\F$ be a field of characteristic $2$, let $V$ be finite-dimensional,
and let
\[
 Q:V\longrightarrow \F,
 \qquad Q(av)=a^2Q(v),
\]
be a quadratic form.  Its polar form is
\[
B(u,v)=Q(u+v)+Q(u)+Q(v).
\]
It is alternating.  Write $R=\rad(B)$, and call $Q$ regular when $R=0$.
In repository terms this is an ordinary characteristic-$2$ $(q,b)$ metric:
the general-bilinear deformation coordinate \texttt{metric.a} is empty.
The Artin--Schreier map and quotient are
\[
 \wpmap:\F\longrightarrow\F,\qquad
 \wpmap(t)=t^2+t,
 \qquad \F/\wpmap(\F).
\]
For a regular form and a symplectic basis $(e_i,f_i)_{i=1}^r$, the usual
characteristic-$2$ Arf class is represented by
\[
 \sum_{i=1}^r Q(e_i)Q(f_i)\pmod{\wpmap(\F)}.
\]
Over a finite field this quotient has two elements and the absolute trace
identifies it with $\mathbb F_2$.  The open question was what replaces that
bit for genuinely transfinite nimber coefficients, where no finite absolute
trace exists.

The answer is: \emph{nothing replaces it over the full nimber field.}  The
class already lives naturally in $\F/\wpmap(\F)$, and this quotient is zero
when $\F=\On$.  More strongly, we give the complete singular normal form
below.

The needed external input is standard.  Conway's nim addition and nim
multiplication turn the class of ordinals into the algebraically closed
characteristic-$2$ field $\On$; Lenstra states this explicitly and analyzes
the algebraic-closure beginning segment~\cite{Lenstra77}.  All remaining
arguments in this note are elementary linear algebra.

\section{Elementary splitting lemmas}

\begin{lemma}[Closure consequences]\label{lem:closure}
Let $\F$ be algebraically closed of characteristic $2$.  Then Frobenius
$x\mapsto x^2$ and the Artin--Schreier map $\wpmap(x)=x^2+x$ are both
surjective.  Hence $\F$ is perfect and $\F/\wpmap(\F)=0$.
\end{lemma}

\begin{proof}
For $a\in\F$, algebraic closure supplies a root of $X^2-a$ and a root of
$X^2+X-a$.  These are respectively a square root of $a$ and a preimage of
$a$ under $\wpmap$.  Frobenius is injective in every field, so the square
root is unique.
\end{proof}

Let $H$ denote the hyperbolic plane: in a basis $e,f$,
\[
 Q(e)=Q(f)=0,\qquad B(e,f)=1.
\]

\begin{lemma}[Every regular plane is hyperbolic]\label{lem:plane}
Let $\F$ be algebraically closed of characteristic $2$.  Every regular
quadratic plane over $\F$ is isometric to $H$.
\end{lemma}

\begin{proof}
Choose a symplectic basis $e,f$, so $B(e,f)=1$, and put
$a=Q(e)$, $b=Q(f)$.  Thus
\[
 Q(xe+yf)=ax^2+xy+by^2.
\]
If $a=0$, then $e$ is a nonzero isotropic vector.  If $a\ne0$, choose
$t\in\F$ satisfying
\[
 t^2+t+ab=0
\]
and set $v=(t/a)e+f$.  Direct substitution gives
\[
 Q(v)=\frac{t^2+t+ab}{a}=0.
\]
In either case take a nonzero isotropic vector $v$.  Regularity supplies
$w_0$ with $B(v,w_0)=1$.  Since $Q(v)=0$,
\[
 w=w_0+Q(w_0)v
\]
satisfies $B(v,w)=1$ and
\[
 Q(w)=Q(w_0)+Q(w_0)B(w_0,v)=0.
\]
Thus $v,w$ are a hyperbolic basis.
\end{proof}

\begin{proposition}[The regular part]\label{prop:regular}
Every regular quadratic space of dimension $2r$ over an algebraically
closed field of characteristic $2$ is isometric to $H^{\perp r}$.
Consequently $\Wq(\F)=0$.
\end{proposition}

\begin{proof}
An alternating nondegenerate form has even dimension and a symplectic
plane.  Lemma~\ref{lem:plane} makes the quadratic form on that plane
hyperbolic.  Its polar-orthogonal complement is again regular.  Induction
splits off $r$ hyperbolic planes.  Since every regular form is metabolic,
its quadratic Witt class is zero.
\end{proof}

The polar radical carries a small residue of information, but it is not an
Arf class.  Let $Z$ denote the one-dimensional zero form and let $A$ denote
the one-dimensional form $x\mapsto x^2$; both have zero polar form.

\begin{lemma}[The radical has only two forms]\label{lem:radical}
Let $R$ be an $s$-dimensional vector space over a perfect field of
characteristic $2$, and suppose the polar form of $Q|_R$ is zero.  Then
there is a unique linear functional $\ell:R\to\F$ such that
\[
 Q(r)=\ell(r)^2.
\]
Up to isometry, $Q|_R$ is $Z^{\perp s}$ if $\ell=0$, and
$A\perp Z^{\perp(s-1)}$ if $\ell\ne0$.
\end{lemma}

\begin{proof}
For a basis $r_1,\dots,r_s$, choose the unique $c_i\in\F$ with
$c_i^2=Q(r_i)$ and define
\[
 \ell\!\left(\sum_i x_i r_i\right)=\sum_i c_i x_i.
\]
Because the polar form vanishes,
\[
 Q\!\left(\sum_i x_i r_i\right)=\sum_i x_i^2Q(r_i)
 =\left(\sum_i c_i x_i\right)^2.
\]
Uniqueness follows from injectivity of Frobenius.  A nonzero linear
functional admits a basis in which it is the first coordinate; this gives
the asserted normal forms.
\end{proof}

\section{Complete normal form over $\On$}

\begin{theorem}[Transfinite nimber quadratic classification]
\label{thm:classification}
Let $Q$ be a finite-dimensional quadratic form over the full nimber field
$\On$, let $B$ be its polar form, and write
\[
 \rank(B)=2r,
 \qquad \dim\rad(B)=s,
 \qquad
 \epsilon=
 \begin{cases}
 0,&Q|_{\rad(B)}=0,\\
 1,&Q|_{\rad(B)}\ne0.
 \end{cases}
\]
Then
\[
 (V,Q)\cong
 \begin{cases}
 H^{\perp r}\perp Z^{\perp s},&\epsilon=0,\\
 H^{\perp r}\perp A\perp Z^{\perp(s-1)},&\epsilon=1.
 \end{cases}
\]
Here the second case necessarily has $s\ge1$.  Thus two such forms are
isometric over $\On$ if and only if the triples $(2r,s,\epsilon)$ agree.
In particular every regular form has a unique isometry class at fixed
dimension, its Arf class is zero, and no transfinite Arf bit exists.
\end{theorem}

\begin{proof}
Choose a vector-space complement $W$ to $R=\rad(B)$.  The restriction
$B|_W$ is nondegenerate: if $w\in W$ is orthogonal to $W$, it is also
orthogonal to $R$, hence lies in $R$, and therefore $w=0$.  Since
$B(W,R)=0$, the quadratic space is the orthogonal sum
\[
 (V,Q)=(W,Q|_W)\perp(R,Q|_R).
\]
Proposition~\ref{prop:regular} identifies the first summand with
$H^{\perp r}$, and Lemma~\ref{lem:radical} gives the two stated radical
normal forms.  Rank, radical dimension, and vanishing of $Q|_R$ are
isometry invariants, so the normal forms are mutually distinct exactly
when their triples differ.
\end{proof}

\begin{remark}[No hidden proper-class step]
The theorem concerns a finite matrix of coefficients and performs finitely
many field operations and root choices.  The fact that $\On$ is a proper
class causes no extra algebraic step: every polynomial used in the proof
has an ordinal root in Conway's field.  Explicitly, start with the set-sized
subfield generated by the finitely many coefficients and take its relative
algebraic closure inside $\On$; this is a set-sized algebraically closed
subfield in which the argument can be carried out.
\end{remark}

The theorem says exactly what survives from the existing
\texttt{ArfInvariants} report after scalar extension to full $\On$:
\[
 \texttt{arf}=0,\qquad
 \texttt{rank}=2r,\qquad
 \texttt{radical\_dim}=s,\qquad
 \texttt{radical\_anisotropic}=\epsilon.
\]
The last field means ``$Q$ is nonzero on the polar radical''; it is not a
remaining anisotropic regular core.

\section{Why the finite Arf bit is still correct}

There is no contradiction between Theorem~\ref{thm:classification} and the
existing finite-nim-field classifier.  Isometry depends on the ground field.

Let $K_d=\mathbb F_{2^d}\subset\On$.  For a regular form over $K_d$, its Arf
class lies in
\[
 K_d/\wpmap(K_d)\cong\mathbb F_2,
\]
where the isomorphism is the absolute trace.  A form of nonzero Arf class
is not hyperbolic over $K_d$, but it becomes hyperbolic after scalar
extension to $K_{2d}$.

\begin{proposition}[Every finite Arf bit dies one stage later]
\label{prop:finite-dies}
The scalar-extension map
\[
 \Wq(K_d)\longrightarrow\Wq(K_{2d})
\]
is zero.  Hence the directed colimit of the finite-field quadratic Witt
groups inside $\On$ is zero.
\end{proposition}

\begin{proof}
For $c\in K_d$, the polynomial $X^2+X+c$ either already has a root in
$K_d$ or is an irreducible quadratic and therefore has a root in
$K_{2d}$.  Thus every $c$ lies in $\wpmap(K_{2d})$, so the induced map
$K_d/\wpmap(K_d)\to K_{2d}/\wpmap(K_{2d})$ is zero.  Arf's classification
then makes the Witt-group map zero.  Every element at every finite stage
therefore maps to zero at a later stage, which is precisely the vanishing
of the directed colimit.
\end{proof}

For the smallest example, the plane $Q(e)=Q(f)=B(e,f)=1$ is anisotropic
over $\mathbb F_2$.  If $t^2+t+1=0$ in $\mathbb F_4$, then $te+f$ is
isotropic, so the same plane is hyperbolic over $\mathbb F_4$.  The finite
Arf bit records failure to split over the declared finite field; it was
never invariant under arbitrary nimber scalar extension.

This settles the proposed ``direct-limit finite-subfield invariant'': the
direct limit is canonically zero.  It also distinguishes two legitimate
questions:
\begin{itemize}
\item Over the ideal full scalar world motivating \texttt{Ordinal}, namely
      $\On$, Theorem~\ref{thm:classification} is the complete classification.
\item If one deliberately freezes the smallest subfield $K$ generated by
      the coefficients and restricts all basis changes to $\mathrm{GL}_n(K)$,
      the relative invariant belongs to $K/\wpmap(K)$.  For transfinite
      coefficients this $K$ need not be finite or algebraically closed.
      This Arf class is generally not by itself a complete higher-dimensional
      isometry invariant over such a field.  That is a different, explicitly
      relative classification problem; it is not a scalar-level invariant of
      $\On$ and has no canonical reduction to one bit.
\end{itemize}

Ogdoad already makes the first distinction for finite metrics: the raw
\texttt{ArfInvariants} result is evaluated at a separately detected degree,
and the higher \texttt{ClassifyWitt}/\texttt{ClassifyBrauerWall} surfaces
store that degree alongside the Arf bit.  The degree is part of those classes,
not decoration.

\section{The Clifford and Brauer--Wall boundary}

The quadratic Witt conclusion immediately answers the Clifford half of the
problem, with one scope qualification.

\begin{corollary}[The regular Clifford class splits]
\label{cor:clifford}
For every regular finite-dimensional quadratic form $Q$ over $\On$, the
graded Morita class of $\Cl(Q)$ on Ogdoad's characteristic-$2$
Clifford/Brauer--Wall surface is zero.
\end{corollary}

\begin{proof}
By Theorem~\ref{thm:classification}, $Q\cong H^{\perp r}$.  Orthogonal sum
maps to graded tensor product of Clifford algebras.  For a hyperbolic basis
$e,f$, the relations
\[
 e^2=f^2=0,\qquad ef+fe=1
\]
identify $\Cl(H)$ with the graded matrix algebra on a $1|1$-dimensional
space, which is graded Morita-trivial.  Therefore
\[
 [\Cl(Q)]=r[\Cl(H)]=0.
\]
\end{proof}

This corollary concerns the subgroup of the graded Brauer--Wall theory
carried by Clifford algebras of regular quadratic forms---the exact surface
used by the repository's finite characteristic-$2$ classifier.  It does not
claim a classification of every possible graded central simple algebra in
characteristic $2$.  For singular $Q$, a nonzero polar-radical vector
becomes a non-scalar homogeneous graded-central Clifford element (the
characteristic-$2$ graded sign is trivial); if its square is zero, it also
generates a nilpotent graded ideal.  Thus the full Clifford algebra is not
graded central simple on the required surface, so there is no Brauer--Wall
class of the whole singular algebra; any polar-nondegenerate complement
nevertheless has the split class above.  Such a complement need not define a
canonical quotient quadratic form when $Q|_{\rad(B)}\ne0$.  The radical
normal-form flag $\epsilon$ is quadratic-form data, not a Brauer--Wall coordinate.

\section{Executable boundary}

The mathematical ground field and the current Rust execution domain are not
the same object.

\begin{itemize}
\item The intended scalar world is full $\On$, but \texttt{Ordinal::nim\_mul}
      is implemented only through a source-verified Kummer window.  Generic
      \texttt{Scalar::mul} is deliberately panic-on-escape.
\item Mathematically, every ordinal below $\omega^{\omega^\omega}$ is
      algebraic over $\mathbb F_2$ and has finite degree.  A finite metric
      there therefore lies in the compositum
      $\mathbb F_{2^m}$, where $m$ is the least common multiple of the entry
      degrees.  The repository's
      \texttt{ordinal\_common\_finite\_subfield\_degree} certifies this field
      when every required Kummer excess is in the verified table and the checked
      degree arithmetic does not overflow, and returns \texttt{None} otherwise.
      Thus there is no mathematical third regime of
      entries in this segment that are not jointly contained in a finite
      nim-subfield, although the partial implementation can honestly report an
      operationally unknown degree.
\item \texttt{arf\_ordinal\_finite} correctly classifies a metric over its
      detected and thereby declared common finite subfield, and uses that
      field's absolute trace.  It does not classify after scalar extension to
      the whole algebraic-closure segment.  Its result is relative to the
      separately detected degree; the higher class wrappers record that degree,
      as
      Proposition~\ref{prop:finite-dies} requires.
\item A metric with coefficients outside the algebraic-closure initial
      segment, or an arbitrary metric over ideal full $\On$, is split over
      $\On$ by Theorem~\ref{thm:classification}; but constructing the isometry
      may require an Artin--Schreier root or square root outside the represented
      multiplication window.  Returning a bare executable ``Arf $=0$'' without
      that distinction would confuse a mathematical base-extension theorem
      with a checked in-window witness.
\end{itemize}

Accordingly no new root-search oracle is required to close the research
problem.  A future implementation may expose a theorem-level
\texttt{classification\_over\_full\_on2} report, or a checked isometry only
when all needed roots are representable.  It should not invent a transfinite
trace to $\mathbb F_2$.

\section{Disposition of \texttt{off}}

The four targets in the former open entry now have exact answers:
\begin{enumerate}
\item \textbf{Domain.}  The ideal full-field classification motivating the
      \texttt{Ordinal} scalar is over $\On$; finite or coefficient-generated
      subfields are explicitly relative refinements.
\item \textbf{Invariant.}  Over full $\On$, the Arf/Artin--Schreier quotient
      vanishes.  The complete possibly singular invariant is
      $(\rank B,\dim\rad B,[Q|_{\rad B}\ne0])$.
\item \textbf{Root route.}  Algebraic closure proves that the required roots
      exist.  The proof gives explicit basis changes once such roots are
      supplied; executable root search remains bounded by the backend rather
      than by mathematics.
\item \textbf{Brauer--Wall.}  Every regular quadratic Clifford class is split;
      singular radicals carry no class of the same kind.
\end{enumerate}

Thus \texttt{off} is resolved by collapse, not by a new transfinite bit.  The
finite bit remains useful precisely because it remembers a smaller declared
ground field.

\begin{thebibliography}{9}

\bibitem{Arf41}
C.~Arf,
\emph{Untersuchungen \"uber quadratische Formen in K\"orpern der
Charakteristik 2. I},
J. Reine Angew. Math. \textbf{183} (1941), 148--167.
\url{https://doi.org/10.1515/crll.1941.183.148}

\bibitem{Conway76}
J.~H. Conway,
\emph{On Numbers and Games},
Academic Press, 1976.

\bibitem{Lenstra77}
H.~W. Lenstra, Jr.,
\emph{On the algebraic closure of two},
Indagationes Mathematicae (Proceedings) \textbf{80} (1977), 389--396.
\url{https://pub.math.leidenuniv.nl/~lenstrahw/PUBLICATIONS/1977e/art.pdf}

\end{thebibliography}

\end{document}
